\documentclass[11pt]{article}

\usepackage[T1]{fontenc}
\usepackage[utf8]{inputenc}
\usepackage{lmodern}
\usepackage[margin=1in]{geometry}
\usepackage{amsmath}
\usepackage{amssymb}
\usepackage{booktabs}
\usepackage{array}
\usepackage{calc}
\usepackage{longtable}
\usepackage{tabularx}
\usepackage{microtype}
\usepackage{xcolor}
\usepackage{fvextra}
\usepackage{seqsplit}
\usepackage{etoolbox}
\usepackage{parskip}
\usepackage[hyphens]{url}
\usepackage[htt]{hyphenat}
\emergencystretch=3em
\sloppy
\setlength{\LTleft}{0pt}\setlength{\LTright}{0pt}
\usepackage[
  colorlinks=true,
  linkcolor=blue!60!black,
  citecolor=blue!60!black,
  urlcolor=blue!50!black,
  pdfcreator={pdflatex},
  pdftitle={Bit-identical transformer inference across ISAs, with a verifiable decode receipt, demonstrated on bare metal},
  pdfauthor={Justin Brian Thompson}
]{hyperref}
\usepackage{bookmark}
\usepackage[round]{natbib}

\DeclareUnicodeCharacter{2212}{\ensuremath{-}}
\DeclareUnicodeCharacter{2264}{\ensuremath{\leq}}
\DeclareUnicodeCharacter{2265}{\ensuremath{\geq}}
\DeclareUnicodeCharacter{2208}{\ensuremath{\in}}
\DeclareUnicodeCharacter{221A}{\ensuremath{\surd}}
\DeclareUnicodeCharacter{2308}{\ensuremath{\lceil}}
\DeclareUnicodeCharacter{2309}{\ensuremath{\rceil}}
\DeclareUnicodeCharacter{230A}{\ensuremath{\lfloor}}
\DeclareUnicodeCharacter{230B}{\ensuremath{\rfloor}}
\DeclareUnicodeCharacter{03A3}{\ensuremath{\Sigma}}
\DeclareUnicodeCharacter{03C3}{\ensuremath{\sigma}}
\DeclareUnicodeCharacter{2032}{\ensuremath{\prime}}
\DeclareUnicodeCharacter{2070}{\textsuperscript{0}}
\DeclareUnicodeCharacter{2075}{\textsuperscript{5}}
\DeclareUnicodeCharacter{2079}{\textsuperscript{9}}


\providecommand{\tightlist}{%
  \setlength{\itemsep}{0pt}\setlength{\parskip}{0pt}}
\newcommand{\pandocbounded}[1]{#1}

\setlength{\emergencystretch}{3em}
\sloppy

\title{\textbf{Bit-identical transformer inference across ISAs, with a verifiable decode receipt, demonstrated on bare metal}}
\author{Justin Brian Thompson \\ Aefinity AI Inc. \\ \texttt{justin.\allowbreak{}brian.\allowbreak{}thompson@gmail.\allowbreak{}com}}
\date{2026-08-27}


\begin{document}

\maketitle

\begin{abstract}
We define CIS-1, a frozen integer-only semantics for the transformer decode path --- every
operation, attention included, in fixed-point integers under one rounding rule --- and show it is
checkable, not merely accurate. Given only the spec text and a public harness, two independent
implementers reproduced the op-level digest on first try {[}A31{]}. The digest is bit-identical across
four x86 microarchitecture/codegen paths {[}A25{]} and on aarch64 {[}A28{]}, and a complete 64-token decode
digest matches across both ISAs {[}A29{]}. A decode receipt --- hashes of the model artifacts, the
tokens, and a hash chain over every step's logit vector --- verifies across the ISA boundary in
public CI {[}A32{]}, and is re-derived bit-for-bit by a UEFI unikernel booting with no OS, on two
physical machines spanning AVX2 and scalar SSE2 {[}A33, A34{]}. The cost is microarchitecture-dependent
and sign-flips: the integer path is 25\% slower than scalar float on Broadwell-U and 4--14\% faster on
Goldmont Plus {[}A26{]}; at parity SIMD, the integer AVX2 kernel is 2.94× faster than the float AVX2
kernel {[}A27{]}. Quality cost, gated
in advance, closed last: the complete all-integer forward pass on a 2-billion-parameter ternary
model costs +0.1239\% perplexity against float --- 40× inside a preregistered +5\% kill line {[}A35{]}. At
that same 2-billion-parameter scale, a full-integer decode receipt minted on x86 verifies
bit-for-bit on aarch64 in public CI, checked by two independent verifiers {[}A39{]}, and the same 2B
receipt is also re-derived by the unikernel booting with no OS, under QEMU {[}A40{]}. The same
unmodified implementation and standalone verifier also handle a second ternary architecture ---
Falcon-E-1B, with a SwiGLU/silu MLP, 16-head/2-KV-head grouped-query attention, a 32,768-token
vocabulary, and untied embeddings --- inside the same +5\% kill line, showing the claim is about the
\end{abstract}

\section{Introduction}\label{introduction}

Ask two machines to run ``the same'' quantized language model on the same prompt and you will
usually get two different sets of logits. The weights are identical; the arithmetic is not.
Parallel reductions sum in different orders, dequantization paths differ between kernels,
normalization layers stay in floating point, and vendor libraries choose algorithms at runtime.
The disagreements are small --- a few units in the last place --- and they are also decisive: greedy
decoding is a chain of argmax operations, and one flipped comparison changes every token that
follows. In practice, ``reproducible'' means ``statistically similar,'' and it is hard to say exactly
what a deployed floating-point model computed on a given input.

This paper makes reproducibility exact and then makes it \emph{checkable}. We define \textbf{CIS-1}, a
canonical integer semantics for transformer inference: a frozen specification in which every
operation of the forward pass --- the ternary matrix--vector product, activation quantization,
RMS normalization, softmax, rotary position embedding, the MLP nonlinearity, and the final argmax ---
is an integer computation with pinned constants and a single normative rounding rule. A
conforming implementation is free to use any instruction set, any vectorization, any
parallelization; it is \emph{not} free to produce a different bit. Conformance is not a matter of
opinion or tolerance: it is a digest. An op-level self-test prints one 64-bit value, and a
complete 64-token greedy decode of a reference model prints another. Either your implementation
prints those values or it does not conform.

Because the semantics are exact, they support something floating-point inference cannot: a
\textbf{decode receipt}. A receipt binds cryptographic hashes of the three model artifacts, the
generated token ids, a digest over those tokens, and a hash chain over the complete integer logit
vector at every decode step. Any conforming machine can replay the decode from source and
verify the receipt bit-for-bit. The receipt is not a claim about a trusted platform; it is a
claim about a computation, and it is checkable by anyone with the artifacts and the spec.

We demonstrate the following, each backed by an entry in the project's public research ledger
with the raw log it was taken from (Table 3):

\begin{enumerate}
\def\labelenumi{\arabic{enumi}.}
\tightlist
\item
  \textbf{The specification is implementable from its text.} Two implementers --- language-model agents
  run under distinct personas, given the specification text and the public conformance harness,
  and denied access to the reference source --- each wrote a from-scratch scalar implementation
  that reproduced the op-level digest on first execution (ledger A31). This is evidence of
  implementability from text, not a third-party audit (§8).
\item
  \textbf{One semantics, two instruction sets.} The op-level digest is identical on four distinct x86
  microarchitecture/codegen paths --- two bare-metal machines spanning AVX2 and SSE2-class hardware,
  a virtualized development host, and full-system emulation (A25) --- and on an aarch64 Neoverse
  system (A28).
  A NEON kernel is bit-identical to the reference on real ARM silicon (A30), and a complete
  greedy decode produces the same token digest on x86-64 and aarch64 (A29). Both digests are
  standing continuous-integration gates: a future divergence fails the build.
\item
  \textbf{The receipt crosses the ISA boundary and the operating-system boundary.} A receipt minted
  on x86 verifies on aarch64 in public CI (A32). The same receipt is re-derived by a unikernel
  that boots from UEFI firmware with no operating system present, on two physical machines ---
  one executing the AVX2 path, the other (which lacks AVX2) the scalar path (A33, A34). At
  BitNet-2B production scale, the same unikernel re-derives the 2B receipt under QEMU/TCG
  emulation --- correctness/identity evidence only, not yet run on physical iron (A40).
\item
  \textbf{Determinism is not the expensive part.} Measured on physical hardware, the integer semantics
  cost 25\% against scalar floating point on a Broadwell-U core and are 4--14\% \emph{faster} on a
  Goldmont Plus core (A26); at parity vector width the integer AVX2 kernel is 2.94× faster than
  the hand-written floating-point AVX2 kernel it replaces (A27). Earlier and larger ``cost of
  determinism'' figures, including this project's own, decompose into the cost of missing SIMD,
  not of integer semantics.
\item
  \textbf{The quality cost is small and was gated in advance.} Against a preregistered +5\% perplexity
  kill line, the all-integer forward pass costs +0.06\% on a 384-hidden reference model (A20), the
  integer-dominant hybrid path (attention still in floating point) costs +0.7408\% on a
  2-billion-parameter ternary model (A21), and the complete all-integer forward pass on the same
  2B model costs +0.1239\% --- closer to float than the hybrid figure, and 40× inside the kill line
  (A35). The first all-integer run at this scale panicked instead: a NORMQ precondition rejected
  an out-of-range residual, tracing to a gap in spec §5.10, where the ACT-I MLP nonlinearity was
  the one op not carrying the spec §5.6 per-vector block exponent that the hybrid boundary already
  applied. The fix --- an RNE-rounded block exponent on the ACT-I output --- is ratified as spec
  erratum v1.0.3; both pinned conformance digests are unchanged.
\item
  \textbf{The token-level identity claim now reaches production scale, across both ISAs.} A complete
  64-token greedy decode of the 2-billion-parameter BitNet-2B model, run in the FullInt
  configuration, prints one digest, reproduced identically across two sequential runs on x86-64:
  \texttt{CIS\_DECODE\ digest=cab11400d737ac4a\ prompt\_toks=4\ gen\_toks=64\ mode=fullint}, and the generated
  text is coherent English (A36). This is the x86 anchor for the BitNet-2B cross-ISA leg --- the 2B
  counterpart of A29 --- and the same decode digest, the decode receipt built on it, and the
  receipt's verification all reproduce bit-for-bit on the GitHub aarch64 runner in public CI (A39).
\item
  \textbf{A third, independent implementation verifies the receipt without the engine.}
  \texttt{cis-verify}, a standalone crate with zero external runtime dependencies and no dependency on
  \texttt{aegis-core}, reproduces both pinned conformance digests and verifies the golden receipt --- all
  six checks (parse, three artifact hashes, prompt tokenization, the 64-step token sequence, the
  cis-digest, and the witness chain) --- in about 1.4 seconds, with tamper tests that fail by naming
  the corrupted field (A37). The verifier itself crosses the ISA boundary: built and run on the
  GitHub \texttt{ubuntu-24.04-arm} runner, the same \texttt{cis-verify} prints
  \texttt{CIS\_SELFTEST\ digest=76985613c965f643\ ALL\_PASS=true}, passes its full suite (81 unit +
  integration tests), and prints \texttt{VERIFY\ PASS} on the x86-minted golden receipt --- both the x86-64
  and aarch64 jobs are now standing CI gates (A38). Honest scope: this was an LM-agent
  transcription of the spec with the reference source visible, not a clean-room audit (§8).
\item
  \textbf{The claim is about the spec, not the checkpoint.} With no engine change and no forge step,
  the unmodified pipeline and the unmodified standalone verifier handle a second ternary
  architecture that is not BitNet --- Falcon-E-1B (24 layers, hidden 2048, SwiGLU/silu MLP,
  16-head/2-KV-head grouped-query attention, 32,768-token vocabulary, untied embeddings):
  full-integer perplexity cost +0.7012\% against float, inside the same +5\% kill line as A35, a
  reproducible 64-token decode digest, and a receipt that verifies bit-for-bit by both the
  reference \texttt{cis\_witness} and the standalone \texttt{cis-verify} (A41).
\end{enumerate}

We are equally explicit about what is \emph{not} shown (§8): the clean-room implementers were
language-model agents rather than third-party engineers; the physical-machine attribution of one
boot log rests on operator witness because the verifier prints no CPU identifier; the
2B-parameter perplexity figures (hybrid and all-integer) use a pruned vocabulary and a short
window and are not comparable to published numbers or to this project's own longer-window
anchors; and no token-level throughput figure for the integer path has yet been
measured. The project publishes negative results and retractions in the same ledger as its
claims, and a standing falsification bounty invites anyone to find a machine on which the digests
diverge.

The contribution, then, is not a faster kernel or a smaller model. It is a way to say, and to
prove, exactly what a model computed --- on hardware you do not have to trust, without an operating
system you would otherwise have to audit.

\begin{center}\rule{0.5\linewidth}{0.5pt}\end{center}

\section{Background and related work}\label{background-and-related-work}

\textbf{Ternary-weight and integer-only LLMs.} BitNet b1.58 quantizes every weight
to \{−1, 0, +1\} (1.58 bits/parameter) and trains directly at that precision
rather than quantizing a float-trained model post hoc; Ma et al.~report
perplexity and end-task parity with FP16 LLaMA baselines at matched size and
training tokens, alongside large memory, latency, and energy reductions
(Ma et al., ``The Era of 1-bit LLMs: All Large Language Models are in 1.58
Bits,'' arXiv:2402.17764 \citep{ma2024bitnet158}, Feb 2024). Microsoft's follow-up technical report
open-sources a 2B-parameter, 4T-token BitNet b1.58 checkpoint with public
weights and CPU/GPU inference code (Ma et al., ``BitNet b1.58 2B4T Technical
Report,'' arXiv:2504.12285 \citep{ma2025bitnet2b4t}, Apr 2025). A companion systems paper describes
CPU kernels for ternary inference and reports CPU speedups of 2.37--6.17×
on x86 and 1.37--5.07× on ARM over unoptimized baselines (Wang et al.,
``1-bit AI Infra: Part 1.1, Fast and Lossless BitNet b1.58 Inference on
CPUs,'' arXiv:2410.16144 \citep{wang2024onebitinfra}, Oct 2024); a related report, bitnet.cpp, targets
edge deployment of ternary models (arXiv:2502.11880 \citep{bitnetcpp2025}). None of this line of
work specifies a bit-exact numerical contract for the non-matmul path
(normalization, softmax, rotary embedding, requantization): the object
released is a model and a fast kernel, not a semantics that a second,
independent implementation is obligated to reproduce bit-for-bit. That gap
--- fast ternary inference exists; a frozen, checkable numerical contract for
it does not --- is the space CIS-1 occupies.

\textbf{Non-determinism and reproducibility in DL inference.} Floating-point
addition is not associative, so reduction order --- which varies with thread
scheduling, tiling, batch size, and kernel choice --- changes results at the
bit level; Shanmugavelu et al.~quantify this across HPC and deep-learning
workloads and evaluate deterministic-summation and hardware-level
mitigations, including PyTorch's determinism flags (arXiv:2408.05148 \citep{shanmugavelu2024nonassoc}, Aug
2024, rev. Oct 2024). Horace He's widely-cited analysis for Thinking
Machines Lab argues the dominant real-world cause of LLM-inference
nondeterminism is not concurrency per se but \emph{batch-size variance}:
individual GPU kernels (matmul, RMSNorm, attention) are themselves
deterministic for a fixed batch shape but change their reduction strategy ---
and therefore their bit-level output --- as batch size changes; the proposed
fix is batch-invariant kernels, measured at roughly 20\% overhead for
matmul/RMSNorm and up to \textasciitilde2× unoptimized for attention (He, ``Defeating
Nondeterminism in LLM Inference,'' Thinking Machines Lab blog, Sep 2025) \citep{he2025defeating}.
DiFR frames the adjacent problem of verifying a provider's inference under
inherent nondeterminism, comparing generated tokens or compressed
activation fingerprints against a trusted reference to catch, e.g.,
undisclosed quantization (Karvonen et al., arXiv:2511.20621 \citep{karvonen2025difr}, Nov 2025).
CIS-1 targets a different point in this space: rather than tolerating
float nondeterminism and verifying \emph{around} it statistically, it removes
floats from the decode path so that reduction order provably cannot change
the result, making the DiFR-style tolerance machinery unnecessary for a
conforming implementation.

\textbf{Integer and fixed-point transformer inference.} I-BERT performs
end-to-end BERT inference in integer-only arithmetic, replacing GELU,
Softmax, and LayerNorm with polynomial/lookup integer approximations (Kim
et al., arXiv:2101.01321 \citep{kim2021ibert}, Jan 2021, rev. Jun 2021). I-ViT extends this to
Vision Transformers with Shiftmax and ShiftGELU, reporting INT8 accuracy
comparable to full precision (Li and Gu, arXiv:2207.01405 \citep{li2022ivit}, Jul 2022, rev.
Aug 2023, ICCV 2023). A 2026 report proposes a clipped-linear softmax
surrogate for int8 edge accelerators (Danopoulos et al., ``Taming the
Exponential,'' arXiv:2604.02292 \citep{danopoulos2026taming}, Apr 2026) --- UNVERIFIED beyond its abstract,
not read in full. These works establish that individual nonlinear ops
\emph{can} be made integer-only; none states its integer approximations as a
frozen, versioned specification with pinned rounding and a conformance
digest, nor addresses RoPE. That distinction --- approximation technique vs.
normative, checkable semantics --- is the same gap noted for BitNet above.

\textbf{Verifiable inference.} Platform (not computational) attestation is
mature: TCG measured boot extends per-stage hashes into TPM PCRs, signed
and reported to a verifier (Trusted Computing Group, ``Overview of TCG
Technologies for Device Identification and Attestation,'' v1.0 rev. 1.37,
Feb 2024) \citep{tcg2024overview}; NIST's SP 800-155 draft specifies BIOS/firmware integrity
measurement guidelines for the same chain (NIST SP 800-155 (Initial Public
Draft), Dec 2011) \citep{nist2011sp800155}. This attests \emph{what code ran}, not what it computed.
Cryptographic approaches instead prove computation: zkLLM proves correct
LLM inference in zero-knowledge with a specialized attention argument,
generating a full-inference proof for a 13B model in under 15 minutes with
sub-200 kB proofs (Sun et al., arXiv:2404.16109 \citep{sun2024zkllm}, Apr 2024, ACM CCS 2024); a
2025--2026 survey catalogs the broader ZKML literature across training,
testing, and inference (Peng et al., arXiv:2502.18535 \citep{peng2025zkmlsurvey}). A parallel line
uses fully homomorphic encryption so a server computes on encrypted inputs
without seeing them, e.g.~Orion for FHE deep learning (arXiv:2311.03470 \citep{orion2023fhe});
FHE and ZK proofs both target privacy/correctness against an untrusted
\emph{executor}, at heavy per-inference cost. CIS-1's decode receipt is a
lighter-weight, complementary object: a commitment (SHA-256 chain) over an
already-computed integer transcript, verified by re-derivation rather than
by circuit proof, and --- because the underlying arithmetic is
order-invariant by construction --- reproducible on different hardware
without the prover/verifier asymmetry ZK and FHE approaches require. It
attests \emph{what was computed}, and is designed to compose with, not replace,
platform attestation.

\textbf{Bare-metal / no-OS LLM inference.} Running LLM inference directly in UEFI
boot-services mode, without an operating system, is not unprecedented: a
secondary report describes a from-scratch C tokenizer/weight-loader/tensor-
math/inference stack running as a UEFI application on a Dell E6510, with no
model, precision, or performance details given beyond an interactive chat
demo and an admission that ``optimization work has barely been done''
(Insights, marvin-42.com, ``Bare-Metal AI: Running LLM Inference Directly in
UEFI, No OS or Kernel Required,'' Mar 2026 \citep{insights2026baremetal} --- itself a secondary report on an
unlinked LocalLLaMA/Reddit project, not a primary technical writeup;
UNVERIFIED beyond this secondary source). That project establishes UEFI as
a viable host for LLM inference; it makes no claim of determinism,
bit-exact reproducibility across ISAs, a cryptographic receipt, or
independent third-party verification --- the properties CIS-1 adds, and the
ones the unikernel demonstration (§5) is built to exercise rather than the
no-OS environment alone.

\textbf{Conformance-by-digest as a design pattern.} Cryptographic standards
have long used exactly this shape: NIST's Cryptographic Algorithm
Validation Program (CAVP) validates a black-box implementation by feeding
it known inputs and checking outputs against known-correct answers (Known
Answer Tests, Multi-block Message Tests, Monte Carlo Tests), most visibly
for AES (NIST CAVP \citep{nistcavp}; AESAVS specification \citep{nistaesavs}, csrc.nist.gov). Conformance is a
match against a pinned test vector, not an audit of source code or
implementation strategy --- implementation freedom plus an exactness
contract, verified by anyone who can run the vectors. CIS-1's two
conformance digests (op-level and token-level) are this pattern applied to
transformer inference: a specification document is normative, and any
implementation --- regardless of ISA, vector width, or language --- either
reproduces the pinned digest or does not conform.

\section{CIS-1: the semantics}\label{cis-1-the-semantics}

CIS-1 rests on one axiom (spec §1): \emph{every reduction in the decode path is a sum of integers whose
worst case provably fits its accumulator.} Integer addition is associative and commutative, so a
conforming implementation may reorder, vectorize, tile, or parallelize any computation --- SSE2,
AVX2, NEON, GPU, FPGA, a for-loop in Python --- and MUST produce bit-identical results. The spec
explicitly rejects the alternative of a canonical float summation order: pinning one loop
structure would kill SIMD freedom and still break under FMA-vs-separate rounding. Determinism is
therefore a property of the arithmetic itself, not a discipline the kernel author must maintain.

Rounding is fixed at one mode everywhere: round-half-even (RNE), with \texttt{rne(num/den)} for \texttt{den\ \textgreater{}\ 0}
computed exactly in \texttt{i128} as the reference primitive (spec §2). The spec calls tie handling normative,
not a detail, because the motivating token flips traced to near-ties decided by float rounding
noise. Four transcendental constants are pinned as the RNE rounding of published 50-digit decimal
expansions, and the integers themselves --- not the decimals --- are normative: \texttt{TWO\_PI\_Q62\ =\ 28976077832308491370}, \texttt{PI\_Q62\ =\ 14488038916154245685}, \texttt{PI\_2\_Q62\ =\ 7244019458077122842}
(independently rounded, not derived from \texttt{PI\_Q62}), and \texttt{LOG2E\_Q32\ =\ 6196328019}. Later pinned
artifacts are verified the same way, by digest rather than by printing the value: the exp LUT
(spec §5.7) carries FNV-1a 64 digest \texttt{0x66C2A0EEB8C2DC43}, golden-tested against an independent
big-integer generator, and the RoPE table for the M7 shape 512×32 at base 10000.0 (spec §5.9) carries
\texttt{0xD8345EBF01E990FA}.

Every object that crosses the decode path is assigned a type, a fixed-point grid, and a proven
bound (spec §3). Quantized activations are symmetric \texttt{i8} in \texttt{−127..+127}, with \texttt{−128} explicitly
forbidden. Dot-product accumulators are \texttt{i32}, valid for \texttt{dim\_in\ ≤\ ⌊i32::MAX/127⌋}. The residual
stream is \texttt{i64} in Q.20 with \texttt{\textbar{}h\textbar{}\ \textless{}\ 2⁵⁰} and vector length \texttt{n\ ≤\ 8192}. q/k/v vectors are \texttt{i32}
Q.16 with \texttt{\textbar{}·\textbar{}\ \textless{}\ 2²⁹} pre-rotation, enforced by a loud panic rather than saturation. Attention
scores are \texttt{i64} Q.24 with \texttt{\textbar{}score\textbar{}\ ≤\ 2⁵⁹}, accumulated in \texttt{i128}. Probabilities live on Q0.15
with \texttt{Σp} within \texttt{⌈T/2⌉} of \texttt{2¹⁵}. exp LUT entries are Q0.31 (spec §5.7), RoPE tables are \texttt{i32} Q1.30
(spec §5.9), norm gains are \texttt{i32} Q.20 in the engine or \texttt{i16} in the reference op, and logits are \texttt{i64}
from an exact integer dot. The spec states its headroom claim plainly: a ternary dot of length L
over \texttt{i8} is bounded by \texttt{127·L}, and a sum of squares by \texttt{127²·L}; violating a stated bound MUST
be a loud failure, never a silent wrap or saturation.

Weight containers pack row-major, four weights per byte, low bit-pair first, with \texttt{dim\_in} a
multiple of 4 (spec §4). The two-bit codes are \texttt{00} = 0, \texttt{01} = +1, \texttt{10} = −1, and --- ratified in this
version --- \texttt{11} decodes to 0 as well (defined-as-zero); a container holding \texttt{11} codes is
conforming, and those positions simply contribute nothing to the dot product.

Section 5 defines twelve operations. \textbf{TMV} (spec §5.1), the ternary matrix--vector product, is exact
in \texttt{i32} in any summation order; equally normative is its rejection surface --- five preconditions
(non-multiple-of-4 \texttt{dim\_in}, an overflowing \texttt{dim\_in}, and three length checks) that every
implementation must reject identically regardless of ISA or internal blocking, with a rejected
call writing no output. \textbf{QUANT-ACT} (spec §5.2) is per-token symmetric absmax quantization onto the
\texttt{i8} grid, \texttt{q\_i\ =\ rne(x\_i·127/absmax)}, with the scale carried forward as the exact rational
\texttt{127/absmax}; the clamp that makes \texttt{−128} unrepresentable is explicit. \textbf{REQUANT} (spec §5.3) applies
a fixed-point multiplier \texttt{y\ =\ clamp(rne((acc·M)/2\^{}(31+S)),\ −127,\ 127)} with \texttt{M\ ∈\ {[}2³⁰,\ 2³¹)} and
\texttt{S\ ≤\ 62}; the \texttt{i64}-input form used by RMSNORM-I is the identical arithmetic widened to \texttt{i128}.
The offline generation of \texttt{(M,S)} from any rational is itself a normative integer procedure,
including the exact \texttt{S}-selection rule and a renormalization case when RNE rounds \texttt{M} up to
\texttt{2³¹}. \textbf{RMSNORM-I} (spec §5.4) specifies every intermediate grid --- sum of squares in \texttt{i64}, an exact
floor integer square root \texttt{t}, an inverse-RMS in Q2.30, and a final Q.20 output through REQUANT ---
because folding these roundings any other mathematically-equivalent way produces different \texttt{i8}
outputs; this full procedure is a v1.0.1 erratum, since the original prose omitted the
bit-determining grids. \textbf{NORMQ} (spec §5.5) is the fused engine form: the RMS factor cancels out of
the \texttt{i8} codes algebraically and survives only in a carried exact rational scale. Spec §5.6 governs
\textbf{container-boundary conversions} --- bf16/f32-to-fixed conversions are exact RNE on the float bit
pattern alone, never on accumulated values, and the block-exponent form \texttt{fix\_f32\_vec} MUST panic
rather than saturate when its fractional width would go negative. Spec §5.7 defines the shared \textbf{exp
machinery}: a parameter-free floor-isqrt chain generates a 1025-entry, strictly-decreasing Q0.31
LUT, looked up by a monotone interpolation. \textbf{SOFTMAX-I} (spec §5.8) works on the Q.24 score grid,
subtracts the max exactly, and normalizes by exact RNE division rather than a Newton reciprocal,
with the declared bound \texttt{\textbar{}Σp\ −\ 2¹⁵\textbar{}\ ≤\ ⌈T/2⌉}. \textbf{ROPE-I} (spec §5.9) tables are not shipped but
generated at load from \texttt{(max\_seq,\ head\_dim,\ f32\ bits\ of\ rope\_theta)} by a normative integer
procedure --- quadrant-reduced sin/cos via 9-term Taylor sums, RNE-requantized to Q1.30 and clamped
to \texttt{±2³⁰}. \textbf{ACT-I} (spec §5.10) elementwise ops work on the Q.20 grid: relu² is an exact squaring
plus one more RNE requant, and silu evaluates one \texttt{exp\_neg} call for \texttt{σ(g)} in Q0.31 before two
further RNE requants. Since the v1.0.3 erratum, the ACT-I output MUST also carry the spec §5.6
per-vector block exponent (\texttt{shift\ =\ max(0,\ bits(max\textbar{}v\textbar{})\ −\ 49)}) before entering NORMQ, identity at
M7 ranges (spec §5.10, ledger A35). \textbf{ARGMAX} (spec §5.11) breaks exact-equality ties on the \texttt{i64} logits to the
lowest index, making a tie a specified, reproducible event rather than rounding noise.

\textbf{Spec §5.12 --- pipeline grid assignments} fix the grid at every stage of the forward pass:

\begin{longtable}[]{@{}
  >{\raggedright\arraybackslash}p{(\columnwidth - 2\tabcolsep) * \real{0.5000}}
  >{\raggedright\arraybackslash}p{(\columnwidth - 2\tabcolsep) * \real{0.5000}}@{}}
\toprule\noalign{}
\begin{minipage}[b]{\linewidth}\raggedright
Stage
\end{minipage} & \begin{minipage}[b]{\linewidth}\raggedright
Grid / accumulator
\end{minipage} \\
\midrule\noalign{}
\endhead
\bottomrule\noalign{}
\endlastfoot
q/k/v (post-rescale) & \texttt{i32}, Q.16 (rotation growth ≤ √2 stays inside \texttt{i32}) \\
Score dot & \texttt{i128} accumulator, scaled by \texttt{1/√head\_dim} in Q0.30, RNE'd to Q.24 \\
V-mix (\texttt{Σ\_t\ p\_t·v\_t}) & exact \texttt{i64}, RNE onto the Q.20 residual grid \\
LM head & exact \texttt{i64} dot, \texttt{i8}-quantized hidden state × Q.20 embedding table \\
\end{longtable}

Section 6 states the exactness contract: an optimized kernel conforms only if it is byte-identical
to the reference for \emph{every} input, not merely inputs a well-behaved caller would produce. The
worked hazard is \texttt{−128}: x86 \texttt{vpsignb} and ARM \texttt{vmulq\_s8} both wrap \texttt{−128\ ×\ −1} within \texttt{i8} where
the reference computes in \texttt{i32}, so both shipped kernels detect \texttt{i8::MIN} in their preparation
pass and route the whole call to the reference --- an unconditional guarantee, not one conditioned
on the caller upholding an invariant. The rejection surface of spec §5.1 is likewise part of the op:
whether a call is rejected must not depend on which CPU the binary landed on. The shipped vector
kernels, \texttt{cis\_avx2} and \texttt{cis\_neon}, are informative rather than normative; their instruction-level
claims are proven by exhaustive enumeration, and whole-kernel equality is checked by test suites ---
deliberately described as a test bar, not a proof.

Section 7 defines what a conforming engine computes: \texttt{CisMode::FullInt} runs embedding lookup, the
residual stream, every norm, every activation quantization, all ternary matvecs, RoPE, attention
scores, softmax, V-mix, the MLP elementwise ops, the LM head, and argmax such that no float value
exists anywhere in the forward pass. Hybrid modes that keep attention or the MLP in \texttt{f32} carry no
cross-ISA claim and sit outside conformance.

Conformance (spec §8) requires all three tiers to pass: Tier 1 op goldens, Tier 2 the operation-level
digest, and Tier 3 the token-level digest. Tier 2 (\texttt{cis\_selftest}) must print exactly:

\begin{Verbatim}[breaklines=true,breakanywhere=true,fontsize=\small]
CIS_SELFTEST digest=76985613c965f643 ALL_PASS=true
\end{Verbatim}

Tier 3 (\texttt{cis\_decode}, prompt \texttt{"Once\ upon\ a\ time"}, 64 new tokens, greedy, EOS ignored, \texttt{FullInt}) must print
exactly:

\begin{Verbatim}[breaklines=true,breakanywhere=true,fontsize=\small]
CIS_DECODE digest=67e8c0a96abc04e1 prompt_toks=4 gen_toks=64 mode=fullint
\end{Verbatim}

\begin{center}\rule{0.5\linewidth}{0.5pt}\end{center}

\textbf{Source map}
- Axiom and rounding/constants paragraphs: spec §1, §2.
- Data types, grids, headroom, and weight packing: spec §3, §4.
- The twelve-operation survey: spec §5.1--§5.12.
- Exactness contract and what a conforming engine computes: spec §6, §7.
- Conformance digests: spec §8.

\begin{center}\rule{0.5\linewidth}{0.5pt}\end{center}

\section{Implementability and cross-ISA identity}\label{implementability-and-cross-isa-identity}

This section is correctness evidence only; no throughput or latency figure appears below
(Rule A --- identity legs carry no performance claim).

\textbf{A four-implementation digest jury (A25).} On 2026-08-02 the op-level self-test printed the
identical digest, \texttt{CIS\_SELFTEST\ digest=76985613c965f643\ ALL\_PASS=true}, on four distinct machines
and codegen paths: the HP Stream (Celeron N4020, Gemini Lake, SSE2-class) running bare iron on
minimal Linux, the Dell i5-5200U running bare iron across two cold boots, the crosvm dev host
(i5-10210U), and QEMU/TCG full-system emulation. The ledger states its own scope exactly, and we
repeat it rather than round it up: ``Four distinct microarchitectures + codegen paths, one
bit-exact integer semantics. Correctness leg only; no performance numbers taken from it'' (A25).

\textbf{The jury crosses its first ISA boundary (A28).} On 2026-08-07 the same op-level digest,
\texttt{76985613c965f643}, reproduced on a GitHub \texttt{ubuntu-24.04-arm} hosted runner (\texttt{uname\ -m}=aarch64
asserted in-job, CPU part \texttt{0xd49}, Neoverse N2 class, Azure southcentralus). The reference
semantics compiled untouched --- \texttt{cis.rs} has zero imports, and \texttt{lib.rs} cfg-gates every
x86-intrinsic module --- and reproduced the digest on first execution. The standing CI gate
\texttt{arm-digest.yml} now pins this digest on every push; a future mismatch fails the build, with the
ledger's own instruction to preserve the log as a finding, not fix it into silence. Because this
leg runs on a cloud runner, the machine is named as precisely as the platform allows; it stands
as an identity artifact only --- no timing is quoted or quotable from it (A28).

Table 2 collects the digest jury; every cell traces to A25 or A28.

\textbf{Table 2 --- the digest jury}

\begin{longtable}[]{@{}
  >{\raggedright\arraybackslash}p{(\columnwidth - 8\tabcolsep) * \real{0.2000}}
  >{\raggedright\arraybackslash}p{(\columnwidth - 8\tabcolsep) * \real{0.2000}}
  >{\raggedright\arraybackslash}p{(\columnwidth - 8\tabcolsep) * \real{0.2000}}
  >{\raggedright\arraybackslash}p{(\columnwidth - 8\tabcolsep) * \real{0.2000}}
  >{\raggedright\arraybackslash}p{(\columnwidth - 8\tabcolsep) * \real{0.2000}}@{}}
\toprule\noalign{}
\begin{minipage}[b]{\linewidth}\raggedright
Machine
\end{minipage} & \begin{minipage}[b]{\linewidth}\raggedright
Microarchitecture
\end{minipage} & \begin{minipage}[b]{\linewidth}\raggedright
Environment
\end{minipage} & \begin{minipage}[b]{\linewidth}\raggedright
Digest
\end{minipage} & \begin{minipage}[b]{\linewidth}\raggedright
Row
\end{minipage} \\
\midrule\noalign{}
\endhead
\bottomrule\noalign{}
\endlastfoot
HP Stream (Celeron N4020) & Gemini Lake, SSE2-class & bare iron, minimal Linux & 76985613c965f643 & A25 \\
Dell Inspiron 15 (i5-5200U) & --- & bare iron (2 cold boots) & 76985613c965f643 & A25 \\
crosvm dev host (i5-10210U) & --- & crosvm (virtualized) & 76985613c965f643 & A25 \\
QEMU/TCG & --- & full-system emulation & 76985613c965f643 & A25 \\
GitHub \texttt{ubuntu-24.04-arm} (Azure southcentralus) & Neoverse N2 class (CPU part 0xd49) & hosted CI runner & 76985613c965f643 & A28 \\
\end{longtable}

\textbf{The specification is implementable from its text (A31).} On 2026-08-08 two independent
implementers --- subagents under distinct personas --- were given the v1.0.1 spec text as their sole
source, with reference-source access forbidden and, as far as the harness can establish, not
used. Each wrote a from-scratch scalar Rust implementation of the spec §5 operations (400 and 484
lines respectively, verified distinct by diff/md5). Both, linked against the verbatim-ported
public conformance harness, printed exactly \texttt{CIS\_SELFTEST\ digest=76985613c965f643\ ALL\_PASS=true}
--- all 14 sections, first run after their own cargo-check. The v1.0.1 spec §5.4 erratum to RMSNORM-I's
Q2.30/2¹⁵ grids is what made this possible: both implementers hit RMSNORM-I's golden on first
attempt from the corrected prose. The ledger states the honest scope plainly, and we do not
soften it: the public conformance suite's golden vectors were visible to the implementers, and
both reports state the A8 goldens were needed to settle \texttt{from\_ratio} edge cases the text leaves
open --- so the proven claim is ``implementable from spec text + public conformance suite, without
reference source,'' not from prose alone. The implementers' ambiguity reports converge on the same
spec gaps, queued as v1.0.2 errata: \texttt{from\_ratio}'s rejection surface and S-selection bracket are
implicit, its RNE-to-2³¹ renormalization at S\textgreater0 is unstated, and QUANT-ACT/RMSNORM-I slice-length
preconditions, empty-ARGMAX, and \texttt{rne\_div} den≤0 behavior are unspecified. Identity artifact only;
dev host (i5-10210U/crosvm) named, no performance figure exists or is claimed (A31).

\textbf{A NEON kernel bit-identical on real ARM silicon (A30).} On 2026-08-07, the same day as A28 and A29 and one day before the clean-room test,
\texttt{cis\_neon::ternary\_matvec\_i8\_neon} --- using \texttt{vmulq\_s8} as the ternary multiply, per-byte bit-pair
extraction through a \texttt{vqtbl1q\_s8} code LUT, exact \texttt{vpaddlq}/\texttt{vpadalq} widening, the shared
precondition check, and the same −128-wrap-to-scalar-fallback contract as the x86 kernel --- was
proven on the GitHub \texttt{ubuntu-24.04-arm} Neoverse runner, all three suites green in one job:
equivalence 5/5 (bit-identity versus the reference, including tails, code \texttt{11}, extremes, and row
independence), contract 6/6 (rejection surface identical to the reference on either side of the
block threshold), and mechanism 3/3 (exhaustive: every byte × 4 pair positions through the LUT,
all 256 activations × 4 codes through \texttt{vmulq\_s8}, and widening-chain extremes). The \texttt{vmulq\_s8}
wrap of \texttt{-128\ ×\ -1} is pinned identical to the x86 \texttt{vpsignb} behavior in both directions by test.
Both standing digest pins (A28's op-level digest and A29's token-level digest) held in the same
run, and the kernel is wired into \texttt{arm-digest.yml} as a standing gate on every push. The ledger is
explicit that no ARM performance number exists or is claimed here --- Rule A requires named
physical hardware, and this runner is shared cloud; the throughput leg stays open until an ARM
board is on the bench (A30).

\textbf{Token-level identity across a complete decode (A29).} \texttt{cis\_decode} runs the full pipeline ---
tokenizer, embeddings, all 7 layers, integer attention, integer LM head, \texttt{argmax\_i64} --- on the
tracked M7 model (hidden 384, 8K vocab) and prints one digest over every token id, prompt
included: \texttt{CIS\_DECODE\ digest=67e8c0a96abc04e1\ prompt\_toks=4\ gen\_toks=64\ mode=fullint}. This
reproduced identically on x86\_64 and aarch64 (GitHub \texttt{ubuntu-24.04} and \texttt{ubuntu-24.04-arm}, both
jobs green in one run), deterministic 3/3 on the dev host; the output is coherent TinyStories
text (``\ldots there was a little girl named Lily\ldots{}'') --- identical \emph{and} fluent. Both digests are standing
CI gates in \texttt{arm-digest.yml}: the x86 job pins the tree, the ARM job pins the ISA, so a future
mismatch is attributable on sight. The pin's design choices are stated in the row and repeated
here rather than left implicit: greedy decoding only, EOS ignored (a fixed 64 tokens, with no
stop-condition dependence), and \texttt{CisMode::FullInt} only --- the Hybrid arms are explicitly
x86-only, because their f32 legs carry no cross-ISA claim. Not established: BitNet-2B scale --- the
2B weights are not in the repository and a 2B decode does not fit in CI; that leg needs a local
ARM board or a paid runner and stays open (A29).

\textbf{Token-level identity reaches production scale, x86 only (A36).} \texttt{cis\_decode} run on the pruned
BitNet-2B artifacts (30 layers, hidden 2560, 50,256-token vocabulary) in the FullInt configuration
prints \texttt{CIS\_DECODE\ digest=cab11400d737ac4a\ prompt\_toks=4\ gen\_toks=64\ mode=fullint}, identical
across two sequential runs on the same host (i5-10210U crosvm). The artifact is the A21/A35 model
with \texttt{\_\_metadata\_\_.aegis\_config} added by \texttt{aegis-forge/add\_aegis\_config.py}; its tensor data blob
is byte-identical to the A21/A35 model, and the same log reproduces the A21/A35 float, hybrid, and
full-integer perplexity figures exactly on this artifact, so the digest and those figures share one
model. The generated text is coherent English (``\ldots in a small town called Greenfield, there lived a
young girl named Lily\ldots{}''). This is the x86 anchor for the BitNet-2B cross-ISA leg --- the 2B
counterpart of A29. Identity evidence only; no timing is quoted or quotable from it (A36).

\textbf{The BitNet-2B cross-ISA leg closes (A39).} The same \texttt{cis\_decode} binary, run on the GitHub
\texttt{ubuntu-24.04-arm} runner (\texttt{uname\ -m}=aarch64 asserted in-job, public run \texttt{33131590730}),
reproduces the identical digest, \texttt{CIS\_DECODE\ digest=cab11400d737ac4a\ prompt\_toks=4\ gen\_toks=64\ mode=fullint}, against the same x86-minted BitNet-2B artifacts; the x86-64 job in the same run
prints the identical line. This upgrades A36 from ``x86 anchor, aarch64 leg open'' to token-level
identity confirmed on both ISAs --- the cross-ISA claim A29 makes for the 7-layer M7 model now
extends to the full 2-billion-parameter model. The digest is a standing CI gate
(\texttt{bitnet2b-receipt.yml}, push to \texttt{main}): a future divergence on either ISA fails the build.
Identity evidence only; no timing is quoted or quotable from it (A39).

\begin{center}\rule{0.5\linewidth}{0.5pt}\end{center}

\section{The decode receipt and bare-metal verification}\label{the-decode-receipt-and-bare-metal-verification}

This section, like §4, carries no timing or throughput number; every figure below is an identity
claim (Rule A).

\textbf{Receipt format (A32).} The golden receipt \texttt{tests/golden/witness\_v1\_m7\_once64.receipt} binds
SHA-256 hashes of all three model artifacts, the 64 generated token ids, the token digest
\texttt{67e8c0a96abc04e1}, and a chained SHA-256 commitment over every decode step's full i64 logit
vector --- chain \texttt{aee25b770bd7b22e…}. It was minted on the crosvm dev host (i5-10210U). This is not
a summary digest: it is the entire integer state sequence, 64 steps deep, each step's complete
logit vector folded into the chain. Every field the receipt binds is listed above exactly as the
ledger row lists it --- the three artifact hashes, the token ids, the token digest, and the
step-by-step logit chain --- because the receipt's value is that a verifier can recompute each of
these independently and compare, not that it asserts a conclusion.

\textbf{Verification crosses the ISA boundary in public CI (A32).} On 2026-08-08, the same receipt
was replayed from source on a GitHub \texttt{ubuntu-24.04-arm} hosted runner, public run \texttt{31249589879},
snapshot \texttt{ce93bbb}: \texttt{artifacts\ 3/3}, the token digest and chain \texttt{aee25b770bd7b22e…} reproduced
exactly, and the harness printed \texttt{VERIFY\ PASS\ —\ replay\ reproduced\ 64\ tokens,\ the\ token\ digest,\ and\ the\ full\ logit\ chain\ bit-for-bit}. The same run's x86-64 job verifies the identical receipt,
and \texttt{arm-digest.yml} now pins both legs on every push. This upgrades the op-level digest's
cross-ISA claim (§4, A28) to a full decode trajectory --- 64 steps times the complete logit vector,
hash-chained --- rather than a single summary value. As with the other cloud-runner legs, the
machine is named as precisely as the platform allows; this is an identity artifact only, and no
timing is quoted or quotable from it (A32).

\textbf{Physical iron, no operating system: the Dell leg (A33).} On 2026-08-08 the Provable AI Kit
stick booted on the Dell Inspiron 15 (i5-5200U, Broadwell-U) with no OS present and re-derived
the golden receipt bit-for-bit. The operator was physically present; boot was via F12 with
Secure Boot off. The firmware appended to BOOTLOG.TXT: \texttt{STAGE\ V:\ witness\ verify\ PASS\ —\ VERIFY\ PASS\ —\ this\ machine\ reproduced\ all\ 64\ decode\ steps\textquotesingle{}\ full\ logit\ vectors\ bit-for-bit,\ with\ no\ OS\ underneath}. The receipt on the stick is md5-identical to the golden (\texttt{87c45bdd…}); the boot
payload is the QEMU-proven \texttt{aegis-kit-iron.img} (\texttt{320e1918…}), and the stick's readback was
verified bit-identical before boot. This completes the chain crosvm-mint → QEMU → public CI
x86-64 → public CI aarch64 (A32) → physical iron, ring 0. We state the attribution limitation as
the ledger states it, not deferred to §8: the new BOOTLOG entry is attributed to the Dell because
it was appended after the baked-in QEMU entry and carries a different firmware memory map
(EMBED.BIN at \texttt{0x3AC000} versus QEMU's \texttt{0x1780000}) --- but verifier mode prints no CPUID, so this
is log-structural evidence, not a hardware identifier (A33).

\textbf{Physical iron, no operating system: the HP leg (A34).} On the same date, the identical stick
and identical receipt were booted on the HP Stream (Celeron N4020, Gemini Lake) for the
first-ever unikernel boot on that machine: VERIFY PASS, payload unchanged since the Dell leg.
BOOTLOG.TXT gained exactly one new entry (diffed against the banked Dell-era log): \texttt{VERIFY\ PASS\ —\ this\ machine\ reproduced\ all\ 64\ decode\ steps\textquotesingle{}\ full\ logit\ vectors\ bit-for-bit,\ with\ no\ OS\ underneath}. The scalar-path corollary matters: the N4020 lacks AVX2, so a PASS on it implies the
SSE2 scalar path executed, since AVX2 would fault with \#UD --- the golden receipt has now been
re-derived bit-for-bit through two disjoint kernel code paths on iron. The attribution limitation
is stated as plainly as the ledger states it: the new entry's firmware memory map is identical to
the Dell entry's, so in-log evidence does not discriminate the two boxes, and attribution of this
entry rests on operator witness (A34). The row also records a firmware finding made incidentally
on this leg: that the N4020's UEFI boots the unikernel at all was previously unknown before this
boot (A34).

\textbf{A third, standalone implementation verifies the receipt without the engine (A37).} Everything
above shows the receipt can be produced and replayed by machines running \texttt{aegis-core} itself, on
either ISA. A receipt that only the reference engine can check is not yet useful to a third party
--- the verifier is what makes it useful, and it must not depend on the code being verified.
\texttt{cis-verify} is a separate crate, written from the spec and the receipt format rather than as a
fork of \texttt{aegis-core}: zero external runtime dependencies, no dependency on \texttt{aegis-core}, no
\texttt{unsafe}, \texttt{no\_std}+alloc at its core. On the same dev host, it reproduces the pinned op-level
digest (\texttt{CIS\_SELFTEST\ digest=76985613c965f643\ ALL\_PASS=true}), both pinned table digests (the exp
LUT and RoPE constants), and the token-level decode digest
(\texttt{CIS\_DECODE\ digest=67e8c0a96abc04e1\ prompt\_toks=4\ gen\_toks=64\ mode=fullint}) on first attempt,
then verifies the golden receipt \texttt{tests/golden/witness\_v1\_m7\_once64.receipt} end to end --- all six
checks (receipt parse, the three artifact hashes, prompt tokenization, the 64-step token-id
sequence, the cis-digest, and the witness chain) --- and prints \texttt{VERIFY\ PASS} in about 1.4 seconds.
Its tamper tests fail by naming the corrupted field (token id, chain, model/vocab hash, or
receipt parse), not by silently passing. Honest scope, stated plainly: this was an LM agent
transcribing the spec with the reference source visible, not a clean-room reimplementation in the
sense A31 uses that term --- it is evidence that the spec and the receipt format are re-implementable
without the engine's SIMD/dispatch code, not an independent third-party audit (§8, A37).

\textbf{Both verifiers cross the ISA boundary, on both receipts (A38, A39).} \texttt{cis-verify}'s standalone
verification (A37) now also runs on aarch64: built and run on the GitHub \texttt{ubuntu-24.04-arm}
runner, it reproduces \texttt{CIS\_SELFTEST\ digest=76985613c965f643\ ALL\_PASS=true}, passes its full suite
(81 unit + integration tests), and prints \texttt{VERIFY\ PASS} on the x86-minted M7 golden receipt
\texttt{tests/golden/witness\_v1\_m7\_once64.receipt}; the x86-64 job in the same run prints the identical
lines. Both are now standing CI gates in \texttt{arm-digest.yml} (A38). The same pattern holds at
BitNet-2B production scale: the receipt minted on x86 (\texttt{tests/golden/witness\_v1\_bitnet2b\_once64.receipt},
cis-digest \texttt{cab11400d737ac4a}, chain
\texttt{917ddf5fea9a848876ddb527d5d5216607637201d6514b94563977009558af32}, bound to artifact
\texttt{facb3597…}) verifies bit-for-bit on the GitHub \texttt{ubuntu-24.04-arm} runner by two independent
implementations at once: the reference \texttt{cis\_witness\ verify} prints \texttt{VERIFY\ PASS}, and the
standalone \texttt{cis-verify} (A37/A38) independently prints \texttt{VERIFY\ PASS} on the same receipt; the
x86-64 job prints the identical lines. This is a standing CI gate, \texttt{bitnet2b-receipt.yml}, on
every push to \texttt{main} (A39). Together, A38 and A39 close the receipt-side counterpart of §4's A39
digest result: the same receipt, at production scale, checked by two independent implementations,
on both ISAs, on every push.

\textbf{The 2B receipt, re-derived by the unikernel with no OS, under QEMU (A40).} On 2026-08-27,
\texttt{aegis-uefi.efi} booted under QEMU/OVMF (TCG, \texttt{-cpu\ max\ -m\ 2048}) from a 1024 MiB FAT32 kit image
carrying the BitNet-2B artifacts --- \texttt{MODEL.SAF} (522,831,917 B), \texttt{EMBED.BIN} (257,310,720 B),
\texttt{VOCAB.BIN} (1,759,936 B) --- and the golden BitNet-2B receipt. BOOTLOG.TXT records \texttt{STAGE\ 2:\ sizes\ OK\ model=522831917\ embed=257310720\ vocab=1759936}, \texttt{STAGE\ 4a}--\texttt{STAGE\ 4d} loading each artifact and
the receipt, \texttt{STAGE\ 5:\ working\ heap\ online}, \texttt{CPUID:\ vendor=AuthenticAMD\ brand="QEMU\ TCG\ CPU\ version\ 2.5+"}, and \texttt{STAGE\ V:\ witness\ verify\ PASS\ —\ VERIFY\ PASS\ —\ this\ machine\ reproduced\ all\ 64\ decode\ steps\textquotesingle{}\ full\ logit\ vectors\ bit-for-bit,\ with\ no\ OS\ underneath}. The serial console shows
\texttt{artifacts:\ 3/3\ hashes\ match}, with receipt and local cis-digest \texttt{cab11400d737ac4a} and chain
\texttt{917ddf5fea9a8488…} agreeing exactly --- the same digest and chain prefix as the x86/aarch64 CI
verifications above (A39). The loader, physical allocator (one contiguous \textasciitilde732 MB claim), and DMA
bounce path handled 782 MB of assets unchanged --- no engine or firmware code change was needed;
only the kit-image packaging script's 64 MiB size constant needed an \texttt{AEGIS\_KIT\_SIZE\_MB} override.
This extends A39 (2B receipt cross-ISA in CI) onto the boot path, and extends the M7-scale iron
result (A33, A34) to BitNet-2B --- but under QEMU/TCG emulation only: correctness/identity evidence
(Rule A), no timing, and no physical-machine claim. A third physical machine (E7c) is staged for
that leg.

\textbf{Where a reviewer should push.} Both physical legs share one structural gap: the verifier
prints no CPU identifier, so the receipt proves what was computed, not unassisted which box
computed it. For the Dell leg, a differing firmware memory map gives log-internal evidence; for
the HP leg, no such internal signal exists, and the claim --- including its AVX2/SSE2 corollary ---
rests on the operator's physical presence at boot (A33, A34). We record this limitation here,
next to the claims that need it, rather than only in the paper's limitations section.

\begin{center}\rule{0.5\linewidth}{0.5pt}\end{center}

\section{Cost and quality}\label{cost-and-quality}

Every figure in this section is a physical-hardware measurement or a same-run,
same-binary comparison logged under \texttt{docs/hardware\_logs/}; none derives from
QEMU/TCG (Rule A), and each number below carries its ledger row and machine
name so it can be checked against the raw log rather than this prose (Rule B).

\subsection{Quality}\label{quality}

The all-integer forward pass costs \textbf{+0.0637\% perplexity} against float on
the M7 reference model (5.643085 vs.~5.639491, pinned 471-token teacher-forced
heldout, i5-10210U crosvm, digest \texttt{0xBED4A17A1A5EE296}, A20) --- comfortably
inside the preregistered +5\% kill line (78× headroom). The earlier
integer-dominant \emph{hybrid} configuration on the same model, in which attention
and activations were still float, cost more: \textbf{+0.3127\%} (5.657126 vs.
5.639491, digest \texttt{0x42E820C2A8A59CD6}, A19). The all-integer path is \emph{closer}
to float than the hybrid, not farther: dropping the two f32→fixed re-entry
quantizations that the hybrid path requires outweighs the Q0.15/Q1.30 table
quantization that the full-integer path adds.

On the production-scale model, BitNet-2B, the measured figure is
\textbf{+0.7408\% perplexity} (integer 30.934140 vs.~float 30.706665,
i5-10210U crosvm, digest \texttt{0x24C4E510A86659D6}, A21). \textbf{This is the HYBRID
path --- attention still runs in f32.} Do not read +0.7408\% as an all-integer
number. The figure also carries caveats that make it non-comparable outside
its own run: it uses a vocabulary pruned from 128,256 to 50,256 tokens
(ASCII-oriented; never compared cross-tokenizer), and a 200-token evaluation
window in an \texttt{\textless{}unk\textgreater{}}-dense region, which is not comparable either to
published full-vocabulary perplexities or to this project's own
longer-window anchor (1,898 tokens, 10.758). The only claim it supports is
the \emph{relative} integer-vs-float cost, measured in the same run with the same
binary and window.

The complete all-integer forward pass on BitNet-2B --- every op, including
attention, in integer --- is now measured: \textbf{+0.1239\% perplexity} (full-integer
30.744724 vs.~float 30.706665, teacher-forced, 199 scored tokens, same
sha-identical artifacts and window as A21, i5-10210U crosvm, argmax digest
\texttt{0xB274DE03F5862DB7}, A35), \emph{closer} to float than the hybrid figure above ---
40× inside the preregistered +5\% kill line. Two sequential runs agreed on
every computed value. Reaching that number required a fix: the unfixed
binary panicked 2/2 runs (\texttt{normq:\ residual\ out\ of\ range}, \texttt{cis\_infer.rs:313}).
An env-gated trace (branch \texttt{cm/e1b-normq-trace}) localized the fault to the
MLP, not attention --- attention re-entry stayed exact at ≤38 bits on all 30
layers, while the ACT-I (relu²) MLP output landed on Q.20 at 52--55 bits, a
genuine spec gap: spec §5.10 lacked the spec §5.6 per-vector block exponent that the
hybrid boundary already carried. The fix --- an RNE-rounded block exponent on
the ACT-I output, degenerating to the identity at M7 ranges --- is ratified as
spec erratum v1.0.3 (spec §11). It changes only that one op: \texttt{CIS\_SELFTEST\ 76985613c965f643} and \texttt{CIS\_DECODE\ 67e8c0a96abc04e1}, the M7 conformance
digests reported throughout this paper, are unchanged. The A21 caveats
(pruned vocabulary, 200-token window, not cross-comparable, no timing) apply
identically to this figure.

\subsection{Generalization to a second model family}\label{generalization-to-a-second-model-family}

Every quality and identity figure above was measured on one architecture family (BitNet/M7's
squared-ReLU MLP, tied embeddings). A41 tests whether CIS-1 is a property of the spec or an
artifact of that one family's numerics, by running the unmodified pipeline, unmodified, on
Falcon-E-1B --- a ternary transformer that is \textbf{not} BitNet: 24 layers, hidden 2048, grouped-query
attention with 16 query heads and 2 KV heads, a SwiGLU/silu MLP (intermediate size 9216) in place
of BitNet's squared-ReLU, untied input/output embeddings, a 32,768-token vocabulary, and
\texttt{rope\_theta} 1e6. The checkpoint is the public \texttt{tiiuae/Falcon-E-1B-Instruct} (TII Falcon-LLM
license); model weights are not redistributed here --- only the receipt, artifact hashes, and logs
are (A41).

No engine code changed and no forge step ran to make this work: the packed artifact's
\texttt{\_\_metadata\_\_.aegis\_config} was already present, and the A35 ACT-I block-exponent fix did not
trigger for this model (A41). Teacher-forced on 198 scored tokens, i5-10210U crosvm: float PPL
\textbf{11.736378}, hybrid-int PPL \textbf{11.823613 (+0.7433\%)}, full-integer PPL \textbf{11.818673 (+0.7012\%)}
--- inside the same preregistered +5\% kill line as A35, and, as with A35, closer to float than the
hybrid path. Hybrid and full-integer argmax digests are \texttt{0x6988214C293A5A8F} and
\texttt{0x4F90C9C906D6A041} respectively (A41). Falcon-E-1B's vocabulary (32,768 tokens, full source
vocabulary) differs from BitNet-2B's pruned tokenizer, so this perplexity is not comparable to the
A19--A21/A35 figures above, or to any other model's perplexity in this paper --- as with every other
quality figure here, only the internal float-vs-integer delta, measured in the same run with the
same binary, is the claim.

The unmodified \texttt{cis\_decode} pipeline prints \texttt{CIS\_DECODE\ digest=3e21adb66d7a17d6\ prompt\_toks=4\ gen\_toks=64\ mode=fullint} on this model, identical across two sequential runs, with the generated
text a coherent English continuation. A receipt, \texttt{tests/golden/witness\_v1\_falcon\_e\_1b\_once64.receipt}
(chain \texttt{3ea8cc522efb28ef…}), mints and verifies bit-for-bit by both the reference \texttt{cis\_witness}
(\texttt{VERIFY\ PASS\ —\ replay\ reproduced\ 64\ tokens,\ the\ token\ digest,\ and\ the\ full\ logit\ chain\ bit-for-bit}) and, independently, the standalone \texttt{cis-verify} (A37/A38) --- all six checks (receipt
parse, the three artifact hashes, prompt tokenization, the 64-step token-id sequence, the
cis-digest, and the witness chain) pass (A41).

The result: the same spec text, the same unmodified engine, and the same standalone verifier
proven against BitNet/M7 above hold, bit-for-bit, on an architecturally distinct ternary model ---
different attention shape, different MLP nonlinearity, different vocabulary, untied embeddings ---
with no engine change and no forge step. This is evidence the claim is about the spec, not one
checkpoint (A41).

\subsection{Cost}\label{cost}

CIS-1's throughput cost was measured directly rather than assumed, on two
physical machines, one binary, three interleaved captures across three
shapes each, 9/9 bit-exact (A26). On the \textbf{Dell i5-5200U (Broadwell-U)},
scalar integer costs 25\% against scalar float: C/B median \textbf{1.248×}
(range 1.239--1.260). On the \textbf{HP N4020 (Gemini Lake / Goldmont Plus)},
integer is \textbf{4--14\% faster}: C/B median \textbf{0.961×} (range 0.856--0.963).
The two machines disagree in sign because the cost is a property of the
microarchitecture, not of the semantics: Broadwell's two FMA units favor the
float LUT+FMA path, while on Goldmont Plus, where floating point is weak,
integer wins. The Dell result decomposes as \textbf{4.61× absent-SIMD × 1.248×
semantics} --- the earlier, larger ``cost of determinism'' figures this project
had previously circulated (including its own) were measuring the cost of
missing SIMD, not of integer semantics. A methodology control on the same
run --- A/B on the HP, where AVX2 is absent and arm A must fall back to arm B,
so the true ratio is known in advance --- measured 1.000× (0.999--1.001, n=9),
which is what licenses trusting the Dell's decomposition.

At parity SIMD width, the picture sharpens further (A27). On the Dell,
\texttt{cis\_avx2::ternary\_matvec\_i8\_avx2} measured \textbf{D/A median 0.340× --- 2.94×
faster} than the hand-written float AVX2 kernel it replaces (range
0.331--0.439), bit-identical to the scalar reference 9/9 with 0 void; D/C
median 0.061× is \textbf{16.4× over the scalar reference}. This result carries a
method note that is part of its provenance, not incidental: the \emph{first}
arm-D run measured D/A 0.276× but was \textbf{rejected} --- arm A's own repeated
measurements drifted +58\% across the run where clock arithmetic predicted
+28\%, while arms B and C tracked the clock to 0.1\%, indicating arm D had
perturbed arm A's cache state and flattered the ratio. The bench was then
changed to re-measure arm A immediately after arm D on every repetition
(A′); on the accepted run, A′/A median = 1.000× (D/A equals D/A′ to three
decimals), which is what licenses treating 0.340× as measured rather than
assumed.

\subsection{Engine context (not CIS-1 claims)}\label{engine-context-not-cis-1-claims}

These figures characterize the surrounding unikernel/kernel-candidate work
on the Dell i5-5200U. They are not part of the CIS-1 conformance or
cross-ISA claims and are reported here only as supporting context.

The ring-0 unikernel was compared against a minimal-Linux decode path under
a preregistered, paired, hands-off-boot protocol (A22): the preregistered
throughput form gave \textbf{+3.6\% / +9.4\% / +5.1\%} across three prompts (3/3),
with 27/27 within-boot bit-exactness and byte-identical responses. The same
protocol's precision prediction, however, \textbf{failed}: \textbf{P-V2-2 FAIL} ---
measured spread 4.7--9.6\% against a predicted \textless3\%. The bit-exactness held;
the wall-time repeatability did not match the prereg's own expectation, and
that failure is reported alongside the pass.

A column-skip kernel candidate measured \textbf{2.88--2.89×} faster than the
incumbent on the real BitNet-2B \texttt{down\_proj} activation distribution (ordered
variant, byte-identical to the incumbent by construction and test --- zero
quality risk on adoption); a chain variant measured 2.80×. Its GMAC/s figures
are stated \textbf{NOMINAL} (skipped work is counted as done; speedup is a time
ratio), and it is not wired end-to-end --- this is a kernel-level result only
(A23).

A memory-bandwidth ceiling was measured on the same machine: peak sequential
read 11.19 / 10.95 GB/s at one thread, 11.63 / 11.70 GB/s at four threads,
against a ternary weight-stream pattern of 0.62 GB/s (A24). The stream
figure is a scalar LUT-walk \textbf{lower bound} on the engine's streaming rate,
not the engine's actual rate --- the bench's own caveat --- and its roughly 18×
gap to peak sequential bandwidth is what motivates the column-skip work
above.

\subsection{Table 4 --- every number in this section}\label{table-4-every-number-in-this-section}

\begin{longtable}[]{@{}
  >{\raggedright\arraybackslash}p{(\columnwidth - 6\tabcolsep) * \real{0.2500}}
  >{\raggedright\arraybackslash}p{(\columnwidth - 6\tabcolsep) * \real{0.2500}}
  >{\raggedright\arraybackslash}p{(\columnwidth - 6\tabcolsep) * \real{0.2500}}
  >{\raggedright\arraybackslash}p{(\columnwidth - 6\tabcolsep) * \real{0.2500}}@{}}
\toprule\noalign{}
\begin{minipage}[b]{\linewidth}\raggedright
Number
\end{minipage} & \begin{minipage}[b]{\linewidth}\raggedright
Metric
\end{minipage} & \begin{minipage}[b]{\linewidth}\raggedright
Machine
\end{minipage} & \begin{minipage}[b]{\linewidth}\raggedright
Row
\end{minipage} \\
\midrule\noalign{}
\endhead
\bottomrule\noalign{}
\endlastfoot
+0.3127\% & Hybrid-path PPL cost, M7 (5.657126 vs 5.639491) & i5-10210U crosvm & A19 \\
+0.0637\% & Full-integer PPL cost, M7 (5.643085 vs 5.639491) & i5-10210U crosvm & A20 \\
+0.7408\% & \textbf{HYBRID}-path (f32 attention) PPL cost, BitNet-2B (30.934140 vs 30.706665) & i5-10210U crosvm & A21 \\
+0.1239\% & \textbf{FULL-INTEGER}-path PPL cost, BitNet-2B, after v1.0.3 erratum (30.744724 vs 30.706665) & i5-10210U crosvm & A35 \\
+0.7433\% & HYBRID-path PPL cost, Falcon-E-1B (11.823613 vs 11.736378) & i5-10210U crosvm & A41 \\
+0.7012\% & FULL-INTEGER-path PPL cost, Falcon-E-1B (11.818673 vs 11.736378) & i5-10210U crosvm & A41 \\
1.248× & Scalar integer/float throughput ratio, C/B & Dell i5-5200U (Broadwell-U) & A26 \\
0.961× & Scalar integer/float throughput ratio, C/B & HP N4020 (Gemini Lake) & A26 \\
4.61× × 1.248× & Absent-SIMD × semantics decomposition & Dell i5-5200U & A26 \\
0.340× (2.94×) & AVX2 integer/float kernel ratio, D/A, parity SIMD & Dell i5-5200U & A27 \\
0.061× (16.4×) & AVX2 integer kernel vs.~scalar reference, D/C & Dell i5-5200U & A27 \\
+3.6\% / +9.4\% / +5.1\% & Ring-0 vs.~minimal-Linux decode throughput, preregistered form & Dell i5-5200U & A22 \\
4.7--9.6\% (predicted \textless3\%, FAIL) & P-V2-2 timing-precision spread & Dell i5-5200U & A22 \\
2.88--2.89× / 2.80× & Column-skip kernel vs.~incumbent, ordered/chain (NOMINAL GMAC/s) & Dell i5-5200U & A23 \\
11.19--11.70 GB/s vs.~0.62 GB/s & Peak sequential read vs.~ternary weight-stream (LUT-walk lower bound) & Dell i5-5200U & A24 \\
\end{longtable}

\begin{center}\rule{0.5\linewidth}{0.5pt}\end{center}

\section{Why this matters}\label{why-this-matters}

\textbf{For deployment without trust.} Edge, air-gapped, and sovereign deployments run models on
hardware and firmware the operator cannot fully audit and networks they cannot rely on. Today the
operator's assurance that ``the model we validated is the model that ran'' rests on platform
attestation --- a measured boot chain that vouches for the software stack --- plus the hope that the
stack is deterministic. CIS-1 removes the hope. A conforming engine's output is a mathematical
function of the artifacts and the input; the decode receipt is the evidence, and any conforming
machine, including one booted from a USB stick with no operating system, can check it. Platform
attestation and computational verification are complementary: one says which binary ran, the
other says what that binary computed.

\textbf{For audit and regulation.} Regulators and incident investigators increasingly need to answer
``what did the system output for this input, and can you prove it?'' With floating-point inference
the honest answer is a distribution. With CIS-1 it is a receipt: hashes of the exact artifacts,
the exact tokens, and a commitment to every intermediate logit vector, replayable by the auditor
on their own hardware. The receipt format is small, the verifier is the reference implementation,
and the specification is public.

\textbf{For research reproducibility.} Every result in this paper is either a digest reproduced in
public continuous integration or a measurement logged on a named physical machine, and Table 3 says
which. A reader who obtains the artifacts does not have to trust our
numbers; they can print the same 64-bit values. We believe this is the correct standard for
inference-engine claims and that the industry's tolerance-based ``matches within epsilon''
reporting has hidden real divergence for years.

\textbf{For the economics of edge AI.} The engine that produces these receipts boots from firmware with
no operating system on decade-old laptops and re-derives the receipt for the reference model there
(A33, A34); the same engine runs a 2-billion-parameter ternary model in its Linux harness on the same
class of hardware (A21). The complete all-integer forward pass on that 2B model now passes its
preregistered quality gate (+0.1239\% perplexity against float, 40× inside the +5\% kill line, A35) ---
though the 2B model has not yet been booted on the unikernel itself; no ledger row yet establishes
that. The cost of the integer semantics is a property of the
microarchitecture, not of the semantics: 25\% against scalar floating point on one core, 4--14\% faster
on another (A26), and at parity vector width the integer AVX2 kernel is 2.94× faster than the
floating-point AVX2 kernel it replaces (A27). On the commodity CPUs where it was measured,
verifiability did not require specialized or trusted hardware, and did not have to cost performance.

\textbf{A standing invitation.} The project maintains a public falsification bounty: find any machine
on which a conforming build fails to reproduce the digests and the author will pay and record the
finding in the research ledger as a deliverable. We would rather buy a counterexample than defend
a claim. Reviewers are invited to try before publication.

\begin{center}\rule{0.5\linewidth}{0.5pt}\end{center}

\section{Limitations and what is not claimed}\label{limitations-and-what-is-not-claimed}

We state these as precisely as the ledger does, because each is where a careful reviewer should push.

\textbf{Scope of the semantics.} The frozen v1.0 spec fixes the data types and grids (spec §3) and lists what is not claimed (spec §10); the numeric operating envelope below is the one the headroom derivations were carried out for, recorded in ledger rows A20--A21 and the v0.3 working notes (CIS-1\_SPEC\_DRAFT\_v0.1.md), not in the frozen text itself. CIS-1 v1.0 covers the BitNet-b1.58-style decoder block
(ternary linear layers, RMSNorm, GQA attention with RoPE, squared-ReLU MLP) with hidden size
384--2560, head dimension ≤128, and sequence length ≤512; headroom bounds are derived for these
ranges only. Bit-identity across kernel paths is claimed for the full-integer pipeline (v0.3 and
later); the earlier hybrid path, which kept attention and activations in floating point, is not
claimed to be cross-path identical and is reported only for the quality comparison in §6.

\textbf{The clean-room implementers were language-model agents.} The two from-scratch implementations
that reproduced the op-level digest (A31) were written by LLM subagents given the specification
text as their sole source, with reference-source access forbidden and, as far as the harness can
establish, not used. This demonstrates implementability from text; it is not an independent
third-party audit, and we do not describe it as one.

\textbf{The \texttt{cis-verify} reimplementation is an LM-agent transcription, not a clean-room audit.} Like
the two implementers behind A31, \texttt{cis-verify} (§5) was built by a language-model agent --- given the
specification and the receipt format as reference material, in three phases --- not by an
independent third party working blind to the source. This demonstrates that the spec and receipt
format are re-implementable without \texttt{aegis-core}'s SIMD/dispatch code and data structures; it is
not the same claim as an external auditor's clean-room verification, and we do not describe it as
one (A37).

\textbf{Attribution of the bare-metal boot logs.} The unikernel's verifier prints no CPU identifier.
The Dell log entry is attributable to that machine because its firmware memory map differs from
the emulator's; the HP log entry's memory map is identical to the Dell's, so in-log evidence does
not discriminate the two machines and attribution of that entry rests on the operator's physical
presence (A33, A34). The scalar-path corollary --- the HP processor lacks AVX2, so a pass on it
implies the SSE2 path executed --- depends on the same attribution.

\textbf{Perplexity figures are not cross-comparable.} The 2B-parameter integer-vs-float comparisons ---
the hybrid path (A21) and the complete all-integer path after the v1.0.3 erratum (A35) --- use a
vocabulary pruned from 128,256 to 50,256 tokens and a 200-token evaluation window in an
out-of-vocabulary-dense region; the resulting absolute perplexities must not be compared with
published full-vocabulary numbers, with this project's own longer-window anchors, or with each
other across runs beyond what the shared, sha-identical artifacts and window support. Only the
\emph{relative} integer-vs-float cost, measured in the same run with the same binary and window, is the
claim.

\textbf{Cost figures are microarchitecture-specific.} The 25\% cost on Broadwell-U and the 4--14\%
advantage on Goldmont Plus (A26) are measurements of two machines, not a model of CPUs in
general; the 2.94× AVX2 result (A27) is one machine, three shapes, three captures, with its
rejected first run and the order control recorded. No token-level throughput number for the
integer path exists yet, and the README says so.

\textbf{Cloud-runner legs are identity evidence only.} The aarch64 results (A28, A29, A30, A32) were
obtained on hosted CI runners; the machine is named as precisely as the platform allows, and no
timing is quoted or quotable from them.

\textbf{The 2B kit has been verified under QEMU but not yet on a physical machine.} A38 and A39
closed the CI cross-ISA gap the previous draft of this paper flagged here, and A40 closes part of
the boot-path gap that remained: the BitNet-2B decode receipt now re-derives bit-for-bit inside the
UEFI unikernel, with no operating system present, under QEMU/TCG emulation. That is correctness and
identity evidence only (Rule A) --- no timing, and not a physical-machine result. What remains open
is the boundary A33 and A34 already crossed, but only at M7 scale: those two legs verified the
smaller M7 receipt on physical Dell and HP hardware, not BitNet-2B. The BitNet-2B kit has not yet
booted on physical iron; a third physical machine (E7c) is staged for that leg.

\textbf{Two model families is still two.} A41 shows the pipeline is not overfit to BitNet's specific
shapes --- it runs unmodified on Falcon-E-1B, a different attention shape (grouped-query, 16h/2kv),
a different MLP nonlinearity (SwiGLU/silu vs.~squared-ReLU), untied embeddings, and a different
vocabulary --- but it is one additional architecture, not a systematic sweep, and both families are
BitNet-style ternary transformers within the class the frozen v1.0 spec targets (spec §3). The
Falcon-E-1B perplexity figures inherit the same non-comparability caveat as A21/A35: its
32,768-token vocabulary differs from BitNet-2B's pruned tokenizer, so the A41 PPL numbers are not
comparable to any other model's PPL in this paper --- only the internal float-vs-integer delta,
measured in the same run with the same binary, is the claim. A harder, qualitatively different
next target is a group-scaled GGUF ternary model --- Ternary-Bonsai-27B --- whose quantization scheme
is not the per-tensor/per-block grid this spec's v1.0 §3 assumes; that generalization has not been
attempted.

\textbf{What the receipt does not do.} A receipt proves that a conforming computation over the bound
artifacts produced the bound outputs. It does not prove which physical machine ran it (that is
platform attestation's job), does not hide the model or the prompt (it is not a zero-knowledge
proof), and does not protect against an adversary who controls the artifacts and the verifier
together.

\textbf{Retractions.} This project has previously published figures it later retracted --- an unlogged
throughput number and two kernel-comparison claims among them. They remain in the ledger, marked,
with the reasons. Every figure in this paper carries a ledger row and a raw log path (Table 3),
and the reader should treat any number without one as a description, not a measurement.

\begin{center}\rule{0.5\linewidth}{0.5pt}\end{center}

\section{Reproducibility statement}\label{reproducibility-statement}

\textbf{Repository.} \texttt{https://github.com/Aefinity-AI/alice-aegis}. The spec is
\texttt{docs/CIS-1\_SPEC\_v1.0.md}; the research ledger with every measurement's raw
log path is \texttt{program/RESEARCH\_LEDGER.md}.

\textbf{CI workflows.} \texttt{.github/workflows/arm-digest.yml} re-proves both
conformance digests and the decode receipt on x86-64 and aarch64 on every
push --- the standing gate behind every identity claim in §4--5. This is a
correctness/identity gate only (Rule A): no timing figure is ever recorded
from it. \texttt{.github/workflows/aefinity-ci.yml} is the general build/format/lint
gate (host job) plus an OVMF boot-correctness job for the unikernel.

\textbf{Golden receipt.} \texttt{tests/golden/witness\_v1\_m7\_once64.receipt}, minted on
the i5-10210U crosvm dev host and verified bit-for-bit on aarch64 in public
CI (run 31249589879, snapshot \texttt{ce93bbb}, A32).

\textbf{The two conformance digests.}
- Op-level: \texttt{CIS\_SELFTEST\ digest=76985613c965f643\ ALL\_PASS=true}
- Token-level decode: \texttt{CIS\_DECODE\ digest=67e8c0a96abc04e1\ prompt\_toks=4\ gen\_toks=64\ mode=fullint}

\textbf{Commands} (the invocation cores of \texttt{.github/workflows/arm-digest.yml}; the workflow additionally pipes each output through \texttt{grep\ \textquotesingle{}\^{}CIS\_…\textquotesingle{}\ \textbar{}\ tee\ /dev/stderr\ \textbar{}\ grep\ -q\ \textless{}digest\textgreater{}\ \textbar{}\textbar{}\ exit\ 1} so a mismatch fails the job, and re-declares the artifact directory \texttt{M} before the witness step --- see the file for the exact wrappers --- which
runs them on every push on both \texttt{ubuntu-24.04} and \texttt{ubuntu-24.04-arm}):

\begin{Verbatim}[breaklines=true,breakanywhere=true,fontsize=\small]
cargo build --release --example cis_selftest --manifest-path aegis-linux/Cargo.toml
./aegis-linux/target/release/examples/cis_selftest

cargo build --release --example cis_decode --manifest-path aegis-linux/Cargo.toml
M=model-lab/tinybit/m7_final_gate_work/artifacts
./aegis-linux/target/release/examples/cis_decode "$M/MODEL.SAF" "$M/EMBED.BIN" "$M/VOCAB.BIN" 64 "Once upon a time"

cargo build --release --example cis_witness --manifest-path aegis-linux/Cargo.toml
./aegis-linux/target/release/examples/cis_witness verify \
  "$M/MODEL.SAF" "$M/EMBED.BIN" "$M/VOCAB.BIN" \
  tests/golden/witness_v1_m7_once64.receipt
\end{Verbatim}

The first two lines reproduce the op-level digest; the next two, the
token-level digest against the in-repo M7 model; the last two replay and
verify the x86-minted golden receipt. A mismatch in any digest is a
falsification of the corresponding claim in §4--5, not a bug to be quietly
fixed --- the workflow says so at the point it would fail.

\textbf{Standalone verifier (\texttt{cis-verify}).} A separate crate at \texttt{cis-verify/}, with zero external
runtime dependencies and no dependency on \texttt{aegis-core} (§5, A37):

\begin{Verbatim}[breaklines=true,breakanywhere=true,fontsize=\small]
cargo test --features std --manifest-path cis-verify/Cargo.toml -- --nocapture

cargo run --release --features std --manifest-path cis-verify/Cargo.toml --bin cis-verify -- \
  tests/golden/witness_v1_m7_once64.receipt \
  "$M/MODEL.SAF" "$M/EMBED.BIN" "$M/VOCAB.BIN"
\end{Verbatim}

The first command runs the crate's test suite, including the pinned op-level digest, both pinned
table digests, the token-level decode digest, and the golden-receipt verification and tamper
tests. The second runs the CLI directly against the golden receipt and prints \texttt{VERIFY\ PASS} (or
\texttt{VERIFY\ FAIL\ (\textless{}field\textgreater{})}) in about 1.4 seconds.

\textbf{BitNet-2B cross-ISA CI (\texttt{bitnet2b-receipt.yml}, A38, A39).} The BitNet-2B artifacts (\textasciitilde745 MB
combined) are too large to commit; they are attached as release assets on tag
\texttt{artifacts-bitnet2b-2026-08-27} and downloaded fresh by every job with
\texttt{gh\ release\ download\ artifacts-bitnet2b-2026-08-27}, then hash-checked against the pinned
SHA-256s before use. \texttt{.github/workflows/bitnet2b-receipt.yml} runs on every push to \texttt{main}, on
both \texttt{ubuntu-24.04} and \texttt{ubuntu-24.04-arm}:

\begin{Verbatim}[breaklines=true,breakanywhere=true,fontsize=\small]
cargo build --release --example cis_decode  --manifest-path aegis-linux/Cargo.toml
cargo build --release --example cis_witness --manifest-path aegis-linux/Cargo.toml
M=2b-artifacts
./aegis-linux/target/release/examples/cis_decode "$M/aegis_pruned_model.cis.safetensors" "$M/embed.bin" "$M/vocab.bin" 64 "Once upon a time"

./aegis-linux/target/release/examples/cis_witness verify \
  "$M/aegis_pruned_model.cis.safetensors" "$M/embed.bin" "$M/vocab.bin" \
  tests/golden/witness_v1_bitnet2b_once64.receipt

cargo run --release --features std --manifest-path cis-verify/Cargo.toml --bin cis-verify -- \
  tests/golden/witness_v1_bitnet2b_once64.receipt \
  "$M/aegis_pruned_model.cis.safetensors" "$M/embed.bin" "$M/vocab.bin"
\end{Verbatim}

The first two lines build the reference drivers; \texttt{cis\_decode} reproduces the BitNet-2B token-level
digest (\texttt{cab11400d737ac4a}); \texttt{cis\_witness\ verify} and the standalone \texttt{cis-verify} CLI each
independently verify the same golden receipt. The aarch64 job running the same four commands is
the source for A38 (the standalone verifier's aarch64 leg) and A39 (the receipt's cross-ISA
verification at 2B scale); the x86-64 job in the same run prints the identical lines.

\textbf{Append-only evidence (Rule C).} \texttt{tests/golden/} and \texttt{docs/hardware\_logs/}
are append-only: every figure in this paper traces to a file under one of
these two paths, and no existing file under either is ever edited, only
added to. Both directories are shipped with the paper's release artifact so
a reader can check a claim against its raw log without re-running anything,
and can re-run the commands above to check it against live hardware anyway.

\begin{center}\rule{0.5\linewidth}{0.5pt}\end{center}

\section{Table 3 --- provenance}\label{table-3-provenance}

Generated by \texttt{docs/paper/gen\_table3.py} from \texttt{program/RESEARCH\_LEDGER.md} rows
A19-A41 and \texttt{docs/CIS1\_PAPER\_OUTLINE.md} §4-6. Do not hand-edit this file ---
edit \texttt{CLAIM\_MAP} in the generator script and re-run it.

Regenerate with:

\begin{Verbatim}[breaklines=true,breakanywhere=true,fontsize=\small]
python3 docs/paper/gen_table3.py
\end{Verbatim}

23 claims total, 23 with a verified existing primary log file.

\begingroup\footnotesize
\begin{longtable}[]{@{}
  >{\raggedright\arraybackslash}p{(\columnwidth - 12\tabcolsep) * \real{0.1429}}
  >{\raggedright\arraybackslash}p{(\columnwidth - 12\tabcolsep) * \real{0.1429}}
  >{\raggedright\arraybackslash}p{(\columnwidth - 12\tabcolsep) * \real{0.1429}}
  >{\raggedright\arraybackslash}p{(\columnwidth - 12\tabcolsep) * \real{0.1429}}
  >{\raggedright\arraybackslash}p{(\columnwidth - 12\tabcolsep) * \real{0.1429}}
  >{\raggedright\arraybackslash}p{(\columnwidth - 12\tabcolsep) * \real{0.1429}}
  >{\raggedright\arraybackslash}p{(\columnwidth - 12\tabcolsep) * \real{0.1429}}@{}}
\toprule\noalign{}
\begin{minipage}[b]{\linewidth}\raggedright
Paper §
\end{minipage} & \begin{minipage}[b]{\linewidth}\raggedright
Claim (short)
\end{minipage} & \begin{minipage}[b]{\linewidth}\raggedright
Value
\end{minipage} & \begin{minipage}[b]{\linewidth}\raggedright
Ledger row
\end{minipage} & \begin{minipage}[b]{\linewidth}\raggedright
Machine
\end{minipage} & \begin{minipage}[b]{\linewidth}\raggedright
Provenance (path or CI run)
\end{minipage} & \begin{minipage}[b]{\linewidth}\raggedright
File exists?
\end{minipage} \\
\midrule\noalign{}
\endhead
\bottomrule\noalign{}
\endlastfoot
4 & Four-implementation digest jury (HP N4020 bare iron, Dell i5-5200U bare iron x2, crosvm, QEMU/TCG) & digest \seqsplit{76985613c965f643}, ALL\_PASS=true, 4 codegen paths & A25 & HP N4020 + Dell i5-5200U + crosvm i5-10210U + QEMU/TCG & \texttt{docs/\allowbreak{}hardware\_\allowbreak{}logs/\allowbreak{}hp\_\allowbreak{}L\_\allowbreak{}BOOTLOG\_\allowbreak{}2026-\allowbreak{}08-\allowbreak{}02.txt} & yes \\
4 & Digest jury crosses to a second ISA (aarch64 CI) & same digest \seqsplit{76985613c965f643} on aarch64 & A28 & GitHub ubuntu-24.04-arm (Neoverse N2) & \texttt{docs/\allowbreak{}hardware\_\allowbreak{}logs/\allowbreak{}cis\_\allowbreak{}selftest\_\allowbreak{}aarch64\_\allowbreak{}github\_\allowbreak{}arm\_\allowbreak{}2026-\allowbreak{}08-\allowbreak{}07.log} & yes \\
4 & Clean-room spec-only reimplementation reproduces the digest & 2 implementers, 400/484 lines, distinct by md5, first-run pass & A31 & i5-10210U crosvm (dev host) & \texttt{docs/\allowbreak{}hardware\_\allowbreak{}logs/\allowbreak{}cis\_\allowbreak{}cleanroom\_\allowbreak{}tier2\_\allowbreak{}2026-\allowbreak{}08-\allowbreak{}08.log} & yes \\
4 & NEON kernel bit-identical on real ARM silicon & equivalence 5/5, contract 6/6, mechanism 3/3 exhaustive & A30 & GitHub ubuntu-24.04-arm (Neoverse N2) & \texttt{docs/\allowbreak{}hardware\_\allowbreak{}logs/\allowbreak{}cis\_\allowbreak{}neon\_\allowbreak{}tests\_\allowbreak{}aarch64\_\allowbreak{}github\_\allowbreak{}arm\_\allowbreak{}2026-\allowbreak{}08-\allowbreak{}07.log} & yes \\
4 & Token-level full-pipeline decode digest identical x86\_64 vs aarch64 & digest \seqsplit{67e8c0a96abc04e1}, prompt\_toks=4 gen\_toks=64 & A29 & GitHub ubuntu-24.04 + ubuntu-24.04-arm & \texttt{docs/\allowbreak{}hardware\_\allowbreak{}logs/\allowbreak{}cis\_\allowbreak{}decode\_\allowbreak{}token\_\allowbreak{}crossisa\_\allowbreak{}ci\_\allowbreak{}2026-\allowbreak{}08-\allowbreak{}07.log} & yes \\
5 & Decode receipt format + CI replay verification & chain \seqsplit{aee25b770bd7b22e}\ldots, CI run \seqsplit{31249589879} (snapshot ce93bbb) & A32 & i5-10210U crosvm (mint) + GitHub ubuntu-24.04-arm + ubuntu-24.04 (verify) & \texttt{docs/\allowbreak{}hardware\_\allowbreak{}logs/\allowbreak{}witness\_\allowbreak{}receipt\_\allowbreak{}aarch64\_\allowbreak{}ci\_\allowbreak{}2026-\allowbreak{}08-\allowbreak{}08.log\ (CI\ run\ \seqsplit{31249589879})} & yes \\
5 & Physical iron verification, Dell (AVX2 path) & STAGE V VERIFY PASS, receipt md5 87c45bdd\ldots, BOOTLOG md5 becd7cef & A33 & Dell Inspiron 15 (i5-5200U, physical iron) & \texttt{docs/\allowbreak{}hardware\_\allowbreak{}logs/\allowbreak{}dell\_\allowbreak{}i5-\allowbreak{}5200U\_\allowbreak{}kit\_\allowbreak{}iron\_\allowbreak{}verify\_\allowbreak{}bootlog\_\allowbreak{}2026-\allowbreak{}08-\allowbreak{}08.txt} & yes \\
5 & Physical iron verification, HP N4020 (SSE2 scalar path) & VERIFY PASS, BOOTLOG md5 4fb5fc8b, one new BOOTLOG entry vs Dell & A34 & HP Stream (Celeron N4020, physical iron) & \texttt{docs/\allowbreak{}hardware\_\allowbreak{}logs/\allowbreak{}hp\_\allowbreak{}n4020\_\allowbreak{}kit\_\allowbreak{}iron\_\allowbreak{}verify\_\allowbreak{}bootlog\_\allowbreak{}2026-\allowbreak{}08-\allowbreak{}08.txt} & yes \\
6 & Quality cost, hybrid path, M7 & +0.3127\% PPL (float 5.639491 vs int 5.657126), digest \seqsplit{0x42E820C2A8A59CD6} & A19 & i5-10210U crosvm & \texttt{docs/\allowbreak{}hardware\_\allowbreak{}logs/\allowbreak{}cis1\_\allowbreak{}e2\_\allowbreak{}int\_\allowbreak{}vs\_\allowbreak{}float\_\allowbreak{}ppl\_\allowbreak{}m7\_\allowbreak{}i5-\allowbreak{}10210U\_\allowbreak{}crosvm\_\allowbreak{}2026-\allowbreak{}08-\allowbreak{}01.log} & yes \\
6 & Quality cost, full-integer path, M7 & +0.0637\% PPL (5.643085 vs 5.639491), digest \seqsplit{0xBED4A17A1A5EE296} & A20 & i5-10210U crosvm & \texttt{docs/\allowbreak{}hardware\_\allowbreak{}logs/\allowbreak{}cis1\_\allowbreak{}fullint\_\allowbreak{}attention\_\allowbreak{}ppl\_\allowbreak{}m7\_\allowbreak{}i5-\allowbreak{}10210U\_\allowbreak{}crosvm\_\allowbreak{}2026-\allowbreak{}08-\allowbreak{}01.log} & yes \\
6 & Quality cost, integer-dominant HYBRID path (f32 attention), BitNet-2B & +0.7408\% PPL (30.934140 vs 30.706665), digest \seqsplit{0x24C4E510A86659D6} & A21 & i5-10210U crosvm & \texttt{docs/\allowbreak{}hardware\_\allowbreak{}logs/\allowbreak{}cis1\_\allowbreak{}e2\_\allowbreak{}bitnet2b\_\allowbreak{}int\_\allowbreak{}vs\_\allowbreak{}float\_\allowbreak{}ppl\_\allowbreak{}i5-\allowbreak{}10210U\_\allowbreak{}crosvm\_\allowbreak{}2026-\allowbreak{}08-\allowbreak{}01.log} & yes \\
6 & Quality cost, FULL-INTEGER path, BitNet-2B (after v1.0.3 erratum) & +0.1239\% PPL (30.744724 vs 30.706665), argmax digest \seqsplit{0xB274DE03F5862DB7} & A35 & i5-10210U crosvm & \texttt{docs/\allowbreak{}hardware\_\allowbreak{}logs/\allowbreak{}cis1\_\allowbreak{}fullint\_\allowbreak{}ppl\_\allowbreak{}bitnet2b\_\allowbreak{}i5-\allowbreak{}10210U\_\allowbreak{}crosvm\_\allowbreak{}2026-\allowbreak{}08-\allowbreak{}27\_\allowbreak{}e1c\_\allowbreak{}fixed.log} & yes \\
6 & Throughput cost by microarchitecture (C/B ratio) & Dell (Broadwell-U) 1.248x slower; HP (Gemini Lake) 0.961x (faster) & A26 & Dell i5-5200U (Broadwell-U) + HP N4020 (Gemini Lake) & \texttt{docs/\allowbreak{}hardware\_\allowbreak{}logs/\allowbreak{}cis\_\allowbreak{}vs\_\allowbreak{}float\_\allowbreak{}L\_\allowbreak{}dell\_\allowbreak{}BOOTLOG\_\allowbreak{}2026-\allowbreak{}08-\allowbreak{}05.txt} & yes \\
6 & AVX2 integer kernel vs float AVX2 kernel & D/A 0.340x = 2.94x faster; D/C 0.061x = 16.4x over scalar & A27 & Dell i5-5200U & \texttt{docs/\allowbreak{}hardware\_\allowbreak{}logs/\allowbreak{}cis\_\allowbreak{}avx2\_\allowbreak{}armD\_\allowbreak{}ordercontrol\_\allowbreak{}L\_\allowbreak{}dell\_\allowbreak{}BOOTLOG\_\allowbreak{}2026-\allowbreak{}08-\allowbreak{}06.txt} & yes \\
6 & Ring-0 unikernel vs minimal Linux decode throughput & +3.6\% / +9.4\% / +5.1\% (prereg form); P-V2-2 FAIL spread 4.7-9.6\% & A22 & Dell i5-5200U & \texttt{docs/\allowbreak{}hardware\_\allowbreak{}logs/\allowbreak{}mech2\_\allowbreak{}U\_\allowbreak{}BOOTLOG\_\allowbreak{}2026-\allowbreak{}08-\allowbreak{}01.txt} & yes \\
6 & Column-skip kernel candidate (engine context, not a CIS-1 claim) & 2.88-2.89x vs incumbent (ordered variant); 2.80x (chain variant) & A23 & Dell i5-5200U & \texttt{docs/\allowbreak{}hardware\_\allowbreak{}logs/\allowbreak{}mech2colskip\_\allowbreak{}L\_\allowbreak{}dell\_\allowbreak{}BOOTLOG\_\allowbreak{}2026-\allowbreak{}08-\allowbreak{}01.txt} & yes \\
6 & Bandwidth ceiling (engine context, not a CIS-1 claim) & peak seq. 11.19/10.95 GB/s (1T), 11.63/11.70 GB/s (4T) vs ternary stream 0.62 GB/s & A24 & Dell i5-5200U & \texttt{docs/\allowbreak{}hardware\_\allowbreak{}logs/\allowbreak{}mech2colskip\_\allowbreak{}L\_\allowbreak{}dell\_\allowbreak{}BOOTLOG\_\allowbreak{}2026-\allowbreak{}08-\allowbreak{}01.txt} & yes \\
4 & Token-level FULL-INTEGER decode digest, BitNet-2B, x86-64 leg & digest \seqsplit{cab11400d737ac4a}, prompt\_toks=4 gen\_toks=64, identical on 2 runs, coherent text & A36 & i5-10210U crosvm & \texttt{docs/\allowbreak{}hardware\_\allowbreak{}logs/\allowbreak{}cis\_\allowbreak{}decode\_\allowbreak{}bitnet2b\_\allowbreak{}fullint\_\allowbreak{}x86\_\allowbreak{}i5-\allowbreak{}10210U\_\allowbreak{}crosvm\_\allowbreak{}2026-\allowbreak{}08-\allowbreak{}27.log} & yes \\
5 & Standalone third-party verifier (no engine dependency) reproduces both digests and verifies the golden receipt & CIS\_SELFTEST \seqsplit{76985613c965f643} ALL\_PASS=true; CIS\_DECODE \seqsplit{67e8c0a96abc04e1}; VERIFY PASS in \textasciitilde1.4s, tamper tests name the field & A37 & i5-10210U crosvm & \texttt{docs/\allowbreak{}hardware\_\allowbreak{}logs/\allowbreak{}cis\_\allowbreak{}verify\_\allowbreak{}standalone\_\allowbreak{}tier2\_\allowbreak{}tier3\_\allowbreak{}golden\_\allowbreak{}i5-\allowbreak{}10210U\_\allowbreak{}crosvm\_\allowbreak{}2026-\allowbreak{}08-\allowbreak{}27.log} & yes \\
5 & Standalone verifier crosses the ISA boundary in public CI & CIS\_SELFTEST \seqsplit{76985613c965f643} ALL\_PASS=true, 81 unit+integration tests pass, VERIFY PASS on x86-minted golden receipt, on aarch64 & A38 & GitHub ubuntu-24.04-arm (Neoverse N2) + ubuntu-24.04 & \texttt{docs/\allowbreak{}hardware\_\allowbreak{}logs/\allowbreak{}cis\_\allowbreak{}verify\_\allowbreak{}aarch64\_\allowbreak{}github\_\allowbreak{}arm\_\allowbreak{}ci\_\allowbreak{}2026-\allowbreak{}08-\allowbreak{}27.log} & yes \\
4 & BitNet-2B decode receipt crosses the ISA boundary in public CI & digest \seqsplit{cab11400d737ac4a} reproduced on aarch64; cis\_witness and standalone cis-verify both print VERIFY PASS & A39 & GitHub ubuntu-24.04-arm (Neoverse N2) + ubuntu-24.04 & \texttt{docs/\allowbreak{}hardware\_\allowbreak{}logs/\allowbreak{}cis\_\allowbreak{}decode\_\allowbreak{}bitnet2b\_\allowbreak{}receipt\_\allowbreak{}crossisa\_\allowbreak{}github\_\allowbreak{}ci\_\allowbreak{}2026-\allowbreak{}08-\allowbreak{}27.log} & yes \\
5 & BitNet-2B receipt re-derived by the unikernel under QEMU/TCG, no OS present & STAGE V VERIFY PASS, cis-digest \seqsplit{cab11400d737ac4a} chain \seqsplit{917ddf5fea9a8488}\ldots, artifacts 3/3 hashes match & A40 & QEMU/TCG (crosvm dev host, i5-10210U) & \texttt{docs/\allowbreak{}hardware\_\allowbreak{}logs/\allowbreak{}unikernel\_\allowbreak{}bitnet2b\_\allowbreak{}verify\_\allowbreak{}qemu\_\allowbreak{}tcg\_\allowbreak{}2026-\allowbreak{}08-\allowbreak{}27.log} & yes \\
6 & CIS-1 generalizes to a second model family (Falcon-E-1B): quality, decode digest, and receipt & full-int +0.7012\% PPL (11.818673 vs 11.736378, 198 tokens); decode digest \seqsplit{3e21adb66d7a17d6}; receipt chain \seqsplit{3ea8cc522efb28ef}\ldots, VERIFY PASS by cis\_witness and standalone cis-verify & A41 & i5-10210U crosvm & \texttt{docs/\allowbreak{}hardware\_\allowbreak{}logs/\allowbreak{}cis\_\allowbreak{}falcon\_\allowbreak{}e\_\allowbreak{}1b\_\allowbreak{}ppl\_\allowbreak{}decode\_\allowbreak{}receipt\_\allowbreak{}i5-\allowbreak{}10210U\_\allowbreak{}crosvm\_\allowbreak{}2026-\allowbreak{}08-\allowbreak{}28.log} & yes \\
\end{longtable}
\endgroup

\begin{center}\rule{0.5\linewidth}{0.5pt}\end{center}

\section{Acknowledgments and reproducibility}\label{acknowledgments-and-reproducibility}

This paper's every quantitative claim traces to a ledger row in
\texttt{program/RESEARCH\_LEDGER.md} and a raw log under \texttt{docs/hardware\_logs/}
(Table 3); nothing here is asserted without that trail, and the repository
ships both directories so a reader can check a claim against its raw log
without re-running anything. Repository:
\texttt{https://github.com/Aefinity-AI/alice-aegis}. The frozen specification is
\texttt{docs/CIS-1\_SPEC\_v1.0.md}; the reproduction commands and CI workflows are in
§9 above.

\textbf{The \$50 falsification challenge.} The author maintains a standing,
public bounty (\texttt{CHALLENGE.md}): find any machine on which a conforming
build produces a different digest or fails witness verification, and the
author will pay \$50 and record the finding in the research ledger as a
deliverable in its own right, not a bug to be quietly fixed. Reviewers are
invited to try before publication.

\nocite{*}
\bibliographystyle{plainnat}
\bibliography{refs}

\end{document}
